\section{Introduction}

Generative-AI products increasingly face litigation over training data, output reproduction, model behavior, and downstream misuse. For engineering organizations, the operational question is not merely whether a lawsuit exists, but whether a theory of liability maps to a product attribute, whether the supporting authority was valid at the relevant filing date, and which mitigation should be reviewed by counsel. Keyword alerts surface noisy filings, docket trackers rarely map cases to features, and conventional retrieval-augmented generation (RAG)~\citep{lewis2020retrieval} treats semantically similar legal documents as comparable even when later decisions have narrowed or superseded them.

We introduce ALERT (\textbf{Agentic Legal Exposure and Risk Tracking}), a continuous sensing system for AI-litigation risk intelligence. ALERT combines CourtListener/RECAP and legal-news ingestion, a 22-field legal-ops schema, and a goal-conditioned reporting loop, producing structured risk registers with citations, provenance, and mitigation candidates while preserving human accountability for legal conclusions.

The paper's main research object is \emph{Product--Litigation Risk Entailment} (PLRE): deciding whether product attributes entail exposure to a legal-risk theory under time-valid precedent. PLRE differs from standard legal textual entailment because the premise is a product attribute, the hypothesis is a risk theory, and the entailment relation is conditioned on a temporal precedent graph. ALERT's graph borrows the legal intuition behind citators such as Shepard's and KeyCite---authority changes over time---but implements a public-corpus, automatically extracted, time-indexed treatment layer and connects it to product-facing entailment rather than to retrospective case lookup.

ALERT is therefore not a legal-monitoring pipeline with an agentic wrapper. Its reusable AI contribution is a point-in-time retrieval operator that changes the candidate evidence distribution \emph{before} generation, rather than filtering citations after an LLM has already produced a report. This distinction matters because stale but lexically similar precedent can dominate ordinary RAG, while PLRE requires product-to-theory entailment under authority valid at the query date. The agentic layer is used to expose unmet evidence atoms and stop when calibrated coverage is satisfied; the novelty claim is the coupling of this verifiable stopping rule with a temporal authority gate whose form is not legal-specific---Eq.~\eqref{eq:tpg_score} applies to any time-stamped, supersession-bearing authority corpus.

Our contributions are: (1) a supersession-aware temporal retrieval mechanism with validity intervals and treatment-edge weighting exposed as a retrieval-time authority gate; (2) PLRE, an \emph{initial} product-conditioned legal-risk entailment benchmark with non-ALERT baselines and ablations separating temporal validity from flat similarity; (3) a goal-conditioned agentic system whose stopping rule controls evidence coverage, not legal correctness; (4) ALERT-Dataset with per-record label provenance, product-mapping annotation guidance, and a human-adjudicated primary circularity check; and (5) a 30-day deployment showing reproducible legal-ops integration.
